% Paper D, revised draft (2026-10-31), restructured as a regular ITP paper after four referee
% reports. Written from latc-complexity at commit 3b6f523 (the last commit that changes a Lean
% file) and from docs/paper-facts.md, which records every theorem statement, file and line used
% here. Lean listings are copied from the source up to `:=`.
%
% Open before submission (see the outline):
%   * anonymity: Paper A and latc-check are cited in the third person, so the text works for
%     double-blind review; the repository link would go through an anonymous mirror;
%   * artifact: public repository or anonymous mirror, Zenodo DOI, CI build (left for later);
%   * page budget: the ITP 2027 regular-paper limit is not yet known.

\documentclass[a4paper,UKenglish,cleveref,autoref,thm-restate]{lipics-v2021}

\usepackage{listings}
\usepackage{array}
\newcolumntype{L}[1]{>{\raggedright\arraybackslash}p{#1}}

\lstdefinelanguage{lean}{
  morekeywords={theorem,def,abbrev,lemma,structure,where,fun,let,if,then,else},
  sensitive=true,
  morecomment=[l]{--},
  morestring=[b]",
}
\lstset{
  language=lean,
  basicstyle=\ttfamily\small,
  keywordstyle=\bfseries,
  columns=fullflexible,
  keepspaces=true,
  xleftmargin=1em,
  breaklines=true,
  breakindent=2em,
  literate=
    {ℕ}{{$\mathbb{N}$}}1
    {∀}{{$\forall$}}1
    {×}{{$\times$}}1
    {∃}{{$\exists$}}1
    {∈}{{$\in$}}1
    {∉}{{$\notin$}}1
    {∧}{{$\wedge$}}1
    {↔}{{$\leftrightarrow$}}1
    {→}{{$\to$}}1
    {↪}{{$\hookrightarrow$}}1
    {≤}{{$\le$}}1
    {∅}{{$\emptyset$}}1
    {ᶜ}{{$^{c}$}}1
    {·}{{$\cdot$}}1
    {₀}{{$_0$}}1
    {ᵢ}{{$_i$}}1
    {φ}{{$\varphi$}}1
    {α}{{$\alpha$}}1
    {β}{{$\beta$}}1
    {₁}{{$_1$}}1
}

\bibliographystyle{plainurl}

\title{Machine-Checked Complexity of Decision-Table Completeness in Lean}
\titlerunning{Machine-Checked Complexity of Decision-Table Completeness}

% Author details as in Paper A. Anonymise if the call asks for double-blind review.
\author{Carlos A. Russell}{Independent Researcher, Buenos Aires, Argentina}{carlos.a.russell@gmail.com}{https://orcid.org/0009-0003-3741-0230}{}
\authorrunning{C.\,A. Russell}
\Copyright{Carlos A. Russell}

\ccsdesc[500]{Theory of computation~Problems, reductions and completeness}
\ccsdesc[300]{Theory of computation~Logic and verification}

\keywords{decision tables, completeness, coNP-completeness, \#P-completeness, Lean 4,
  Cook--Levin theorem}

\supplement{The Lean development, latc-complexity}
% TODO: archive with a DOI at submission (Zenodo), or an anonymous mirror for review.

\hideLIPIcs
\nolinenumbers

\begin{document}

\maketitle

\begin{abstract}
A decision table is complete when every combination of its declared condition values that the
domain constraints allow is covered by a row. We prove in Lean~4, over multi-tape Turing
machines and with no unproved assumptions beyond the definitions themselves, that for finite-domain decision tables with clause
constraints and rows that are complete condition tuples, deciding whether some combination is
missing is NP-complete, deciding completeness is coNP-complete, and counting the missing
combinations is \#P-complete. An adequacy theorem
ties these word problems to the set-based definitions of an existing formalisation of
decision-table completion, so the results speak about the original objects. Complementary
polynomial-time results locate the hardness: all three problems are polynomial without
constraints, and with a fixed number of attributes, where one machine runs within
$c\,(|x|+1)^{dn+d}$ on instances with $n$ attributes. For monotone two-literal constraints,
counting the missing combinations remains \#P-complete under reductions with a single oracle
call; for this we prove a
one-call version of Valiant's theorem on \#Monotone-2SAT in the same Turing-machine model.
\end{abstract}

%-----------------------------------------------------------------------------------------------
\section{Introduction}
\label{sec:intro}

Decision tables state business rules, regulations and clinical guidelines as rows of
condition values~\cite{PollackHicksHarrison1971}. A table is \emph{complete} when every combination of declared condition
values that the domain constraints allow is covered by some row; the uncovered combinations
form its \emph{missing set}~$M(T_0)$. Completeness checking is a staple of decision-table
practice, from early expert-system checkers~\cite{SuwaScottShortliffe1982} to the analysers of
the DMN standard~\cite{CalvaneseDLMMT2016}, which report missing rules. Counting them matters
too: one missing edge case and a table that misses half of its valid inputs are both
``incomplete'', and $|M(T_0)|$ is what tells them apart.

Russell~\cite{PaperA} re-examined a decision-table completion method over declared finite
domains and argued, on paper, that completeness under clause constraints is coNP-complete and
that counting $|M(T_0)|$ is \#P-complete (Proposition~5 there). The accompanying Lean
development, latc-check~\cite{latc_check_0_1_0}, formalised the decision-table model and a
combinatorial reduction from its own propositional formulas, but not the complexity classes: those need a model of computation,
polynomial-time bounds and encodings, which latc-check did not have.

This paper supplies them. We compose latc-check's model with Complexitylib~\cite{Complexitylib},
a Lean~4 library of complexity theory on multi-tape Turing machines with a machine-checked
Cook--Levin theorem, and prove the classification end to end, together with results that neither
earlier development contains. On paper the classification is elementary: with no rows,
completeness is unsatisfiability of the constraints. The contribution is its mechanisation on
Turing machines with an adequacy theorem, the tractability boundary with an explicit degree, and
the one-call Valiant theorem. Every theorem we cite depends only on Lean's three standard
axioms, and the build contains no \lstinline{sorry}.

\emph{Scope.} Throughout, a row of the table is a complete tuple of condition values, as in
latc-check. Tables with compact rows, such as DMN rules with wildcards or ranges, are a different
problem: there completeness is coNP-hard even without constraints (\Cref{sec:related}).

\subparagraph*{Contributions.}
\Cref{tab:contrib} separates what this paper adds from what it reuses.
\begin{enumerate}
\item \textbf{The classification, machine-checked end to end} (\Cref{sec:decision,sec:counting}):
  non-emptiness of $M(T_0)$ is NP-complete, completeness coNP-complete and counting \#P-complete,
  for finite domains with clause constraints, with every number of the input written in binary.
  The reductions produce only instances with Boolean domains and empty tables, so hardness
  holds already there (by inspection of the reduction, not as a separate Lean theorem).
\item \textbf{Adequacy} (\Cref{sec:encoding}): the encoded problems are exactly latc-check's
  set-based ones.
\item \textbf{One-call hardness for monotone two-literal counting} (\Cref{sec:monotone}): a
  definition of one-call reductions in Complexitylib's model, and a one-call version of Valiant's
  theorem, from which counting the missing combinations stays \#P-complete under one-call
  reductions when every constraint is $A_i = 1 \lor A_j = 1$, for tables with arbitrary rows (hard already with none).
\item \textbf{The tractability boundary} (\Cref{sec:boundary}): polynomial time without
  constraints, and with a fixed number of attributes, for both decision problems and for
  counting, with an explicit degree linear in the number of attributes.
\end{enumerate}

\begin{table}[t]
\caption{What is reused and what is new.}
\label{tab:contrib}
\centering\small
\begin{tabular}{L{0.38\textwidth}L{0.26\textwidth}L{0.26\textwidth}}
\hline
result & before this paper & here \\
\hline
decision-table model, missing set & latc-check (Lean) & reused \\
SAT reduction to missing states & latc-check, from its own formulas & new, from Complexitylib's CNF \\
NP/coNP-completeness, \#P-completeness & on paper~\cite[Prop.~5]{PaperA} & Turing machines, in Lean \\
adequacy of the encoded problems & -- & new \\
polynomial cases: no constraints, fixed $k$ & analogues for firewalls~\cite{HamiltonRoughanNguyen2022} & new for decision tables \\
explicit degree in the number of attributes & -- & new \\
one-call reductions, one-call Valiant & Valiant (Turing reductions); Senellart--Gnatenko (logical model) & new in Complexitylib's model \\
\hline
\end{tabular}
\end{table}

\Cref{sec:prelim} recalls Complexitylib and latc-check, \Cref{sec:encoding} fixes the input
language, \Cref{sec:decision,sec:counting,sec:monotone,sec:boundary} prove the results,
\Cref{sec:related} discusses related work and \Cref{sec:formal} the formalisation.
\Cref{app:statements} lists the main Lean statements and \Cref{app:artifact} the artifact.

%-----------------------------------------------------------------------------------------------
\section{Preliminaries}
\label{sec:prelim}

\subparagraph*{Machines, encodings and classes.}
Complexitylib, built on CSLib~\cite{CSLib2026} and Mathlib, works with multi-tape Turing
machines over a binary alphabet. A type $\alpha$ is
given an encoding, an injection $\alpha \hookrightarrow \{0,1\}^*$; pairs, lists and natural
numbers have standard encodings, numbers in binary (\lstinline{encNat}) or in unary
(\lstinline{encUnary}). A function $f$ is computable in polynomial time between two encodings
when some machine maps the encoding of every $a$ to the encoding of $f(a)$ within a polynomial
in the input length. A language is in P when its characteristic function is; in NP when it is
the set of inputs having a certificate of polynomially bounded length accepted by a
polynomial-time verifier; and coNP is the class of complements. A function $f$ into
$\mathbb{N}$ is in \#P when $f(a)$ is the number of certificates of length exactly $p(|a|)$,
for a polynomial~$p$, that a polynomial-time verifier accepts on~$a$. A library lemma,
\lstinline{mem_SharpP_of_polyTime_length}, lets the certificate length be any polynomial-time
computable function of the input instead; our counting verifiers use it.

\subparagraph*{Reductions.}
Three notions are used. A \emph{many-one} reduction from $A$ to $B$ is a polynomial-time $r$
with $a \in A \iff r(a) \in B$; NP-, coNP-hardness and completeness are defined with it. A
\emph{parsimonious} reduction from $f$ to $g$ is a polynomial-time $r$ with $f(a) = g(r(a))$;
Complexitylib defines \#P-hardness with it. These two are Complexitylib's. The third is defined
in this development; it is the polynomial-time metric reduction attributed to
Krentel~\cite{Krentel1988}:
\begin{definition}[one-call reduction]
\label{def:onecall}
$f$ reduces to $g$ with one call if there are polynomial-time functions $r$ and
$\mathit{post}$ with $f(a) = \mathit{post}(a, g(r(a)))$ for every~$a$, where $\mathit{post}$
reads $a$ and the answer $g(r(a))$ written in binary. $g$ is \emph{one-call \#P-hard} if every
function in \#P, under any encoding, reduces to it with one call, and \emph{one-call
\#P-complete} if it is moreover in \#P.
\end{definition}
The running time of $\mathit{post}$ is polynomial in the length of $a$ and of the oracle's
answer. For a target $g$ whose values have polynomially many bits, as for every function in
\#P, a one-call reduction is a polynomial-time Turing reduction making exactly one oracle query,
whose answer is post-processed together with the input, and one-call hardness implies hardness
under polynomial-time Turing reductions, the notion of Valiant~\cite{Valiant1979enumeration}.
(Without that bound the notion would be too weak: a $g$ with doubly exponential values gives
$\mathit{post}$ exponential time in $|a|$ while carrying no information. We use one-call
hardness only for \#Monotone-2SAT, whose value is at most $2^{\#\text{variables}}$, and for a
target proved to be in \#P.) Parsimonious reductions are one-call, with
$\mathit{post}(a, y) = y$; so \#P-hardness implies one-call \#P-hardness, and a
parsimonious reduction followed by a one-call reduction is one-call.
Zank\'o defines \#P-completeness with a variant of many-one reductions and shows the
$(0,1)$-permanent complete under it~\cite{Zanko1991}. One-call hardness is weaker than parsimonious
hardness, which matters below: a problem whose decision version
is trivial cannot be parsimoniously \#P-hard unless P~=~NP, but can be one-call \#P-hard.
Our hardness proofs start from Complexitylib's Cook--Levin theorem and its parsimonious form:
SAT is NP-hard, UNSAT coNP-complete, and \#SAT \#P-hard under parsimonious reductions.

\subparagraph*{latc-check's model.}
Attributes are indexed by a finite type~$\iota$, attribute $i$ having the declared finite domain
$D_i$, so the space of condition tuples $\mathcal{D} = \prod_i D_i$ is finite. Constraints are a
decidable predicate $C$ on tuples, and a table is represented by its set $T_0$ of condition
tuples. The valid space is $\mathcal{V} = \{x \in \mathcal{D} : C(x)\}$, the missing set is
$M(T_0) = \mathcal{V} \setminus T_0$, and the table is complete when $M(T_0) = \emptyset$. In
Lean these are \lstinline{validSpace C} and \lstinline{missing C T}, both \lstinline{Finset}s, so
duplicate rows play no role; this model, not latc-check's own reduction, is what the present
development reuses. Two lemmas check that latc-check's clause constraints and satisfiability agree
with Complexitylib's CNF formulas (\lstinline{holdsAll_eq_eval},
\lstinline{satisfiable_decode_iff}); the hardness proofs below use a new reduction stated
directly on Complexitylib's CNF.

%-----------------------------------------------------------------------------------------------
\section{Instances, encoding and adequacy}
\label{sec:encoding}

\subparagraph*{Instances.}
An instance is a triple $I = (d, S, T)$ of lists of natural numbers:
\begin{itemize}
\item $d = [d_0, \dots, d_{n-1}]$, the domain sizes; attribute $i$ ranges over
  $\{0, \dots, d_i - 1\}$, and $n$, the number of attributes, is the length of~$d$;
\item $S$, the clause constraints: a list of clauses, each a list of conditions $(i, v)$, read
  $A_i = v$;
\item $T$, the rows of $T_0$: a list of lists of values.
\end{itemize}
A \emph{state} is a list $s$ of length $n$ with $s_i < d_i$. It satisfies the condition $(i, v)$
when $i < n$ and $s_i = v$, and it is \emph{valid} when every clause of $S$ has a condition it
satisfies. Every triple is an instance; nothing is rejected as malformed, and the semantics
settles each degenerate case. A condition on a nonexistent attribute, or with a value outside
the domain, is never satisfied; an empty clause makes every state invalid; a domain of size
zero leaves no states; a row of the wrong length or with an out-of-domain value covers no state;
and repeated rows or conditions change nothing. Then
\begin{align*}
\textit{MissingNonempty} &= \{I : \text{some valid state is not a row of } T\},\\
\textit{Complete} &= \overline{\textit{MissingNonempty}},
\end{align*}
and $\textit{missingCount}(I)$ is the number of valid states that are not rows.

\subparagraph*{A running example.}
Take two attributes with $d = [3, 2]$, the constraint $A_0 = 0 \lor A_1 = 1$, and the rows
$(0,0)$, $(0,1)$, $(1,1)$:
\[
I = \bigl([3, 2],\ [[(0,0), (1,1)]],\ [[0,0], [0,1], [1,1]]\bigr).
\]
Of the six states, four are valid: $(0,0)$, $(0,1)$, $(1,1)$ and $(2,1)$. The first three are
rows, so $M(T_0) = \{(2,1)\}$: the instance is in \textit{MissingNonempty}, not complete, and
$\textit{missingCount}(I) = 1$.

\subparagraph*{Encoding and input size.}
The main encoding, \lstinline{encInstBin}, writes the triple with Complexitylib's encodings of
pairs and lists and every number in binary: the domain sizes, the attribute index and value of
each condition, and the values of each row. The number of attributes is not written separately;
it is the length of the list~$d$. The input size $|x|$ in every polynomial bound below is the
length of this encoding. With binary domain sizes, $|\mathcal{D}| = \prod_i d_i$ can be
exponential in $|x|$ even for one attribute.

Checking a condition $(i, v)$ against a state reads $s_i$ at a binary index~$i$. We use
Complexitylib's \lstinline{listGetD_binary}, which walks the list while reading the bits of the
index and returns a default when the list ends, so the cost is polynomial in the lengths of the
list and of the index's binary encoding, not in the index's value, which is never written in
unary. The values are shifted by
one before the lookup, so that the default~$0$ matches no value: an out-of-range attribute fails
the condition, as the semantics requires. A second encoding, \lstinline{encInst}, writes the
attribute index in unary. The identity is polynomial-time from \lstinline{encInst} to
\lstinline{encInstBin}, so hardness for the unary encoding transfers to the binary one; every
result below holds for both, except the explicit-degree bounds of \Cref{sec:boundary}, which are
stated for \lstinline{encInstBin} only.

\subparagraph*{Adequacy.}
An instance denotes a table in latc-check's model: attributes \lstinline{Fin n}, domains
\lstinline{Fin dᵢ}, the constraint predicate given by the clauses, and as $T_0$ the tuples whose
list of values is a row. The word problems are exactly latc-check's:
\begin{lstlisting}
theorem mem_MissingNonempty_iff (I : Inst) :
    I ∈ MissingNonempty ↔ (missing I.con I.table).Nonempty
theorem mem_Complete_iff (I : Inst) : I ∈ Complete ↔ missing I.con I.table = ∅
\end{lstlisting}
A third theorem, \lstinline{coe_missing_map_toList}, shows that the tuples of $M(T_0)$, read as
lists, are exactly the witnesses of \textit{MissingNonempty}, and \textit{missingCount} is
defined as the cardinality of latc-check's missing set. Every complexity result below is
therefore a statement about latc-check's \lstinline{missing} on clause-constrained tables; it
says nothing about constraint languages other than clauses.

%-----------------------------------------------------------------------------------------------
\section{The decision problems}
\label{sec:decision}

\begin{theorem}
\label{thm:decision}
Under \lstinline{encInstBin}, \textit{MissingNonempty} is NP-complete and \textit{Complete} is
coNP-complete. The reductions use only instances with domains $\{0,1\}$ and $T_0 = \emptyset$.
\end{theorem}

\begin{proof}[Proof sketch]
\emph{Membership.} The certificate is a missing state. The verifier checks that it lies in the
domains, satisfies every clause and is not a row; it runs in polynomial time
(\lstinline{checkBin_polyTime}). A state in the domains has $s_i < d_i$ for every~$i$, so its
encoding is no longer than that of~$d$, which bounds the certificate length by~$|x|$.
\emph{Hardness.} Given a CNF formula $\varphi$, renumber its variables densely as
$0, \dots, n-1$ and build the instance with $n$ attributes of domain $\{0,1\}$, one clause per
clause of $\varphi$ (the literal $x_i$ becomes $A_i = 1$, and $\lnot x_i$ becomes $A_i = 0$), and
no rows. Its valid states are the satisfying assignments of~$\varphi$, and with no rows they are
all missing. So $\varphi$ is satisfiable iff the instance is in \textit{MissingNonempty}, and
unsatisfiable iff it is complete: a reduction from SAT and from UNSAT, polynomial-time
(\lstinline{fromSAT_polyTime}). Composed with Complexitylib's Cook--Levin theorems and with the
identity re-encoding, this gives the theorem (\lstinline{MissingNonempty_npComplete_bin},
\lstinline{Complete_coNPComplete_bin}). By adequacy the same holds for
$\{I : M(T_0) \neq \emptyset\}$ and $\{I : M(T_0) = \emptyset\}$ stated with latc-check's
\lstinline{missing}.
\end{proof}

The empty table shows where the hardness lies: in the constraints, not in the rows. With
$T_0 = \emptyset$, completeness is unsatisfiability of the constraints. Rows that are complete
tuples only remove single states, and \Cref{sec:boundary} shows that without constraints they
cannot make the problems hard; compact rows could (\Cref{sec:related}). The
development also has a separate Boolean model, whose attributes are the variables of a
Complexitylib CNF formula used as the constraints and whose rows are bit vectors; NP-, coNP- and
\#P-completeness are proved for it too (\Cref{app:statements}), without an adequacy theorem.

%-----------------------------------------------------------------------------------------------
\section{Counting}
\label{sec:counting}

\begin{theorem}
\label{thm:counting}
Under \lstinline{encInstBin}, \textit{missingCount} is \#P-complete under parsimonious
reductions; the hardness reduction uses only instances with domains $\{0,1\}$ and
$T_0 = \emptyset$.
\end{theorem}

\begin{proof}[Proof sketch]
\emph{Membership.} Membership needs every missing state to have exactly one accepted
certificate. Let $L$ be the length of the encoding of the domain sizes~$d$. A certificate is a
bit string of length~$L$. The verifier decodes the first $n$ numbers from it into a state~$s$
and accepts when $s$ is missing \emph{and} the certificate equals the canonical certificate of
$s$: the encoding of $s$ followed by zeros up to length~$L$. The encoding of a state in the
domains is no longer than that of $d$, so every missing state has a canonical certificate of
length~$L$; the canonical certificate determines the state; and any other string decoding to $s$
is rejected. So the accepted certificates of length $L$ are in bijection with $M(T_0)$
(\lstinline{missingCount_eq_acceptCount}), and both the verifier and $L$ are polynomial-time.
Without the canonicity test, a state could have several accepted certificates (different
padding) and the count would be wrong.

\emph{Hardness.} The reduction of \Cref{thm:decision} is parsimonious: the satisfying
assignments of $\varphi$ and the missing states of its instance are in bijection, so
$\#\mathrm{SAT}(\varphi) = \textit{missingCount}(\mathit{fromSAT}(\varphi))$
(\lstinline{sharpSAT_eq_missingCount}). Composition with Complexitylib's parsimonious
Cook--Levin theorem gives \lstinline{missingCount_sharpPComplete_bin}.
\end{proof}

In the separate Boolean model there is also a parsimonious reduction \emph{to} \#SAT: add to the
constraints, for each row of the right length, the clause stating that the tuple differs from
that row (\lstinline{countMissing_eq_sharpSAT}).

%-----------------------------------------------------------------------------------------------
\section{Monotone two-literal constraints and Valiant's theorem}
\label{sec:monotone}

Russell~\cite{PaperA} notes that counting the valid states stays \#P-complete when every
clause is $A_i = 1 \lor A_j = 1$ over domains $\{0,1\}$, citing
Valiant~\cite{Valiant1979enumeration}. Deciding such instances is easy: with domains
$\{0,1\}$ and only positive two-literal clauses on existing attributes, the all-ones state is
valid; the shape tests also require every condition to name an existing attribute, so every
such instance with no rows has a missing state (\lstinline{mem_missingNonempty_of_pos2B}). So counting cannot
be \#P-hard under parsimonious reductions unless P~=~NP, and the right notion is the one-call
reduction of \Cref{def:onecall}.

Let \textit{SharpMonotone2SAT} be the number of models of a CNF formula whose clauses are all
$x_i \lor x_j$, and \textit{countValidMonotone2} the number of valid states of an instance with
domains $\{0,1\}$, clauses of two conditions with value~1, and no rows. Let
\textit{countMissingMonotone2} be the same without the condition on rows: the number of missing
combinations $|M(T_0)|$ of an instance with that domains and constraints and any rows
(\lstinline{countMissingMonotone2_eq_card_missing}). As with Complexitylib's \#3SAT, the three
functions are total and return~$0$ on inputs of any other shape; the shape tests are
polynomial-time. Both counts are in \#P, and \#Monotone-2SAT reduces to each
parsimoniously (\lstinline{sharpMonotone2SAT_reduces} and its variant with suffix
\lstinline{_countMissing}): the reduction produces either a monotone instance with no rows, where the two counts agree, or a fixed
instance of the wrong shape, where both are~$0$. What remains is the hardness of \#Monotone-2SAT
itself.

\begin{theorem}[one-call Valiant]
\label{thm:valiant}
\#Monotone-2SAT is one-call \#P-hard. Hence \textit{countValidMonotone2} and
\textit{countMissingMonotone2} are one-call \#P-complete under \lstinline{encInstBin} and
\lstinline{encInst}.
\end{theorem}

Valiant proved \#Monotone-2SAT \#P-complete under polynomial-time Turing reductions,
interpolating over many oracle calls. Senellart and Gnatenko machine-check the one-call form in a
logical model of \#P with first-order reductions~\cite{SenellartGnatenko2026}. We prove it in
Complexitylib's model by one reduction from \#SAT, which is \#P-hard parsimoniously, so by
\Cref{def:onecall} it suffices to reduce \#SAT to \#Monotone-2SAT with one call.

\subparagraph*{Step 1: models as independent sets.}
Let $\varphi$ have variables $0, \dots, n-1$ (after dense renumbering) and clauses
$c_1, \dots, c_m$. Build a graph $G$ with a vertex $(i, b)$ for each variable $i$ and value $b$,
and a vertex $o$ for each literal occurrence, read ``$o$ is the first true literal of its
clause''. Edges join $(i, 0)$ and $(i, 1)$; any two occurrences of the same clause; an
occurrence and the value of its variable that falsifies it; and an occurrence and each value
that makes an \emph{earlier} literal of its clause true. The vertices fall into $K = n + m$
cliques, $\{(i,0),(i,1)\}$ for each variable and the occurrences of each clause, so every
independent set has at most $K$ vertices. An independent set of exactly $K$ vertices picks a
value for each variable, which is an assignment~$\sigma$, and an occurrence in each clause that
$\sigma$ makes true, with every earlier literal false. That occurrence is unique given~$\sigma$,
so these independent sets are in bijection with the models of~$\varphi$
(\lstinline{modelCount_eq_card}). The ``first true literal'' choice is what makes the
correspondence exact, for any CNF with densely numbered variables (the hypotheses of
\lstinline{modelCount_eq_card}): repeated and complementary literals need no special case. An
empty clause has no vertex, so no independent set reaches $K$, and the count is $0$, as it should
be.

For example, $\varphi = x_0 \lor x_1$ has $n = 2$, $m = 1$, $K = 3$, vertices $(0,0)$,
$(0,1)$, $(1,0)$, $(1,1)$, $o_1$ for $x_0$, $o_2$ for $x_1$. The edges at $o_1$ go to $o_2$ and
$(0,0)$; those at $o_2$ go to $o_1$, $(1,0)$ and $(0,1)$. The independent sets of size~3 are
$\{(0,1), (1,0), o_1\}$, $\{(0,1), (1,1), o_1\}$ and $\{(0,0), (1,1), o_2\}$: three, one per model.

\subparagraph*{Step 2: a monotone formula.}
Let $N$ be the number of vertices of $G$, and give every vertex $t = n + 1$ fresh pendant
leaves. The monotone formula $\psi$ has a variable per vertex and leaf, and a clause $x_u \lor
x_w$ per edge $\{u, w\}$ of the graph with leaves. Its models are the vertex covers, the
complements of the independent sets. An independent set of the graph with leaves consists of an
independent set $S$ of $G$ together with any subset of the leaves of the $N - |S|$ vertices
outside~$S$, so
\[
x \;=\; \#\mathrm{SAT}(\psi) \;=\; \sum_{S \text{ independent in } G} 2^{\,t\,(N - |S|)}
\qquad (\text{\lstinline{modelCount_psi}}).
\]
Every vertex has leaves, so no vertex of the formula is isolated, and $\psi$ has exactly two
positive literals per clause.

\subparagraph*{Step 3: one digit.}
Write $x$ in base $2^t$. Multiply by $2^{tK}$ and divide by $2^{tN}$:
\[
\Bigl\lfloor \frac{x \cdot 2^{tK}}{2^{tN}} \Bigr\rfloor
= \Bigl\lfloor \sum_{S} 2^{\,t\,(K - |S|)} \Bigr\rfloor .
\]
Since $|S| \le K$, every exponent is a nonnegative multiple of $t$, so the sum is an integer and
the floor is exact. Sets with $|S| < K$ contribute multiples of $2^t$, and sets with $|S| = K$
contribute $1$ each. Their number is $\#\mathrm{SAT}(\varphi) \le 2^n < 2^t$. So
\[
\#\mathrm{SAT}(\varphi) \;=\; \Bigl\lfloor \frac{x \cdot 2^{tK}}{2^{tN}} \Bigr\rfloor \bmod 2^t
\qquad (\text{\lstinline{sharpSAT_eq_post}}).
\]
Multiplying before dividing avoids computing $N - K$, which can be negative when $\varphi$ has
empty clauses ($K$ counts them, $N$ does not). Then every term $2^{t(K-|S|)}$ has $|S| < K$, so the
quotient is a multiple of $2^t$ and the digit is $0$, which is correct.
The post-processing needs only binary multiplication, division, remainder and powers of two on
numbers of polynomially many bits, all polynomial-time in Complexitylib; the construction of
$\psi$ is polynomial-time; and the oracle is called once, on~$\psi$
(\lstinline{toMonotone_polyTime}, \lstinline{post_polyTime}, \lstinline{sharpSAT_oneCallReducible}).
Composing with the parsimonious \#P-hardness of \#SAT gives
\lstinline{sharpMonotone2SAT_oneCallSharpPHard}, and with the reductions to
\textit{countValidMonotone2} and \textit{countMissingMonotone2}, the second half of
\Cref{thm:valiant} (\Cref{app:statements}). The development also keeps
the conditional form, \lstinline{countValidMonotone2_oneCallSharpPComplete_bin}, which takes
\#Monotone-2SAT's hardness as a hypothesis; \Cref{thm:valiant} discharges it.

%-----------------------------------------------------------------------------------------------
\section{Where the hardness comes from}
\label{sec:boundary}

The reductions of \Cref{sec:decision,sec:counting} use clause constraints and as many
attributes as the formula has variables. Removing either ingredient makes all three problems
polynomial-time (\Cref{tab:boundary}).

\subparagraph*{No constraints.}
Then every state is valid and $|M(T_0)| = \prod_i d_i - r$, where $r$ is the number of
\emph{distinct} rows that are states (\lstinline{missingCount_eq_of_noConstraints}). The product
can be exponential in $|x|$, but its binary encoding is no longer than those of the factors
together, so a fold that multiplies in binary stays polynomial: each intermediate product has
at most $|x|$ bits. The count is therefore polynomial-time, and so are completeness and
non-emptiness, which compare $\prod_i d_i$ with~$r$ (\lstinline{missingCountNoConstraints_polyTime},
\lstinline{completeNoConstraints_mem_P}, \lstinline{missingNonemptyNoConstraints_mem_P}).

\subparagraph*{At most $k$ attributes: deciding.}
With binary domain sizes, $\mathcal{D}$ cannot be enumerated even for $k = 1$. Let $N$ be the
number of value entries written in the instance, in conditions and in rows, and call a number
\emph{unwritten} if it occurs in none of them. An unwritten value satisfies no condition, and a
state containing one is not a row; so in a missing state, an unwritten value can be replaced by
any other unwritten value of the same domain, and the state stays missing. At most $N$ of the
numbers $0, \dots, N$ are written, so one of them, $u$, is unwritten. If a missing state uses an
unwritten value $v > N$ in attribute~$i$, then $d_i > v > u$, so $u$ is in the domain and can
replace~$v$. Hence if a missing state exists, one exists whose values lie among $0, \dots, N$
and the written values: at most $2N+1$ candidates per attribute and $(2N+1)^k$ candidate
states (\lstinline{length_candidateStates_le},
\lstinline{exists_missing_candidate_iff}). Checking each candidate is polynomial, so both decision
problems are in P for each~$k$ (\lstinline{missingNonemptyFixedAttrs_mem_P},
\lstinline{completeFixedAttrs_mem_P}).

\subparagraph*{At most $k$ attributes: counting.}
A candidate state stands for all the states obtained by replacing its unwritten values, so
counting must weigh it. For each attribute, split its domain into classes: each written value
inside the domain alone, and one class holding all the unwritten values of the domain (possibly
empty), represented by~$u$. Which
conditions a state satisfies, and whether it is a row, depend only on its vector of classes, so
\[
|M(T_0)| = \sum_{\substack{c \text{ a class vector,}\\ \text{valid and not a row}}}\
\prod_{i} |c_i|
\qquad (\text{\lstinline{missingCount_eq_weighted_of_le}}),
\]
over at most $(N+1)^k$ class vectors (\lstinline{length_classStates_le}). Each product is at most
$\prod_i d_i$, so the sum is computed in binary by the bounded-growth fold of the
no-constraint case, and the count is polynomial-time for each~$k$
(\lstinline{missingCountFixedAttrs_polyTime}).

\subparagraph*{An explicit degree.}
``Polynomial time for each $k$'' allows the polynomial to depend on $k$ arbitrarily. We bound it
with Complexitylib's \lstinline{ComputableTidilyInO}, which keeps explicit bound functions and a
single constant through composition, case splits and folds. The following are proved, each
for non-emptiness, completeness and counting, each for time and space, under
\lstinline{encInstBin}:
\begin{enumerate}
\item \emph{One machine for all instances.} There are $c, d$ such that one machine decides
  non-emptiness on every instance within $c\,(|x|+1)^{dn+d}$, where $n$ is the instance's number
  of attributes (\lstinline{missingNonempty_explicitDegree}; \Cref{app:statements}).
\item \emph{Uniform in $k$}, a direct corollary of~1 (\lstinline{missingNonemptyFixedAttrs_uniform}).
  There are $c, d_0, d_1$ such that for every $k$, on instances with at most $k$ attributes, one
  machine decides the problem within $c\,(|x|+1)^{d_0 k + d_1}$.
\item \emph{$k$ in the input.} There are $c, d$ such that one machine, given $k$ in unary and an
  instance, decides ``at most $k$ attributes and $M(T_0) \neq \emptyset$'' within
  $c\,(|(k,x)|+1)^{dk+d}$ (\lstinline{missingNonemptyWithBound_explicitDegree}).
\end{enumerate}
The degree grows linearly with the number of attributes. This is an XP bound, $|x|^{O(k)}$, not
fixed-parameter tractability, $f(k) \cdot |x|^{O(1)}$; we claim no FPT result.

\emph{What is not proved.} For each fixed $k$, a separate machine with $k$ built in that rejects
instances with more than $k$ attributes, with constants independent of~$k$. Testing $n \le k$
inside the machine goes through Complexitylib's case-split and constant-output lemmas, whose
constants depend on the machine, and making them uniform would need changes to the library.
Statements~2 and~3 cover the same ground: the first on the instances where the answer can be
positive, the second on all instances with $k$ given as input.

\begin{table}[t]
\caption{Where the hardness comes from.}
\label{tab:boundary}
\centering\small
\begin{tabular}{L{0.11\textwidth}L{0.13\textwidth}L{0.12\textwidth}L{0.15\textwidth}L{0.16\textwidth}L{0.14\textwidth}}
\hline
rows & constraints & attributes & deciding & counting & hardness via \\
\hline
tuples & clauses & unbounded & NP-, coNP-complete & \#P-complete & many-one, parsimonious \\
tuples & none & unbounded & P & polynomial-time & -- \\
tuples & clauses & at most $k$ & P, degree $O(k)$ & polynomial-time, degree $O(k)$ & -- \\
tuples & monotone two-literal & unbounded & trivial if $T_0 = \emptyset$ & one-call \#P-complete & one-call \\
compact (wildcards) & none & unbounded & coNP-complete (on paper) & \#P-complete (on paper) & DNF tautology, \#SAT \\
\hline
\end{tabular}
\end{table}

%-----------------------------------------------------------------------------------------------
\section{Related work}
\label{sec:related}

\subparagraph*{Compact rows.}
Our rows are complete tuples, as in latc-check. In DMN and in many rule bases a row is compact: a
conjunction of tests, with wildcards or ranges, standing for many tuples. Then the rows form a
DNF formula, completeness is its tautology problem, coNP-complete even without
constraints~\cite{Cook1971}, and the uncovered inputs are the models of its negation, a CNF, so
counting them is \#SAT. Russell~\cite{PaperA} makes the same observation. Elmallah and Gouda's
firewall rules are compact rows over integer ranges~\cite{ElmallahGouda2017}, and Thiemann and
Yamada's patterns are compact rows over terms~\cite{ThiemannYamada2026}. With complete tuples as
rows, as here, the rows are not the source of hardness: the constraints are
(\Cref{sec:boundary}).

\subparagraph*{Hardness of coverage.}
Hardness of rule coverage is not new in itself. Elmallah and Gouda prove firewall completeness
NP-hard by reduction from 3-SAT, with no membership result~\cite{ElmallahGouda2017}. Thiemann
and Yamada prove pattern completeness co-NP hard and verify their decision procedure in
Isabelle, the hardness proof being on paper~\cite{ThiemannYamada2026}. Kapur et al.\ showed
quasi-reducibility co-NP complete~\cite{KapurNRZ1991}, and the volume of a union of boxes is
\#P-hard~\cite{BringmannFriedrich2010}. For firewalls with a fixed number of header fields,
Hamilton, Roughan and Nguyen give polynomial-time algorithms~\cite{HamiltonRoughanNguyen2022},
an analogue of our fixed-$k$ results.

\subparagraph*{Decision tables and tabular specifications.}
Checking tables for completeness and consistency is long established: in expert-system
checkers~\cite{SuwaScottShortliffe1982}, for decision tables~\cite{VanthienenMA1998}, for SCR and
RSML requirements~\cite{HeitmeyerJL1996,HeimdahlLeveson1996}, and for DMN, where Calvanese et al.\
give analysis algorithms~\cite{CalvaneseDLMMT2016,CalvaneseDLMMT2018} and, with background
knowledge, an ExpTime upper bound~\cite{CalvaneseMDM2019}. PVS generates coverage and
disjointness obligations for tabular specifications~\cite{OwreRushbyShankar1997}: proof
obligations about given tables, not complexity results. Menghi et al.\ prove completeness of
Simulink Requirements Tables undecidable, by reduction from Hilbert's tenth problem, because
their conditions range over real-valued signals, and check it with
SMT~\cite{MenghiEtAl2025,ZanengaEtAl2025}. Vu and Demetrovics study decision tables whose entries
may be unknown: for one given input, deciding its possible and certain decisions is NP- and
coNP-complete, and counting the valuations yielding a decision \#P-hard, all by encoding a formula
in the constraints on the unknowns; a fixed number of unknowns is polynomial, another analogue of
our fixed-$k$ case~\cite{VuDemetrovics2026}. They define completeness for a single input but state
no complexity for deciding it. Their model asks what a partly unknown table says about one input;
ours asks whether a fully specified table covers every input its declared domains and constraints
allow.

\subparagraph*{Rule-base verification.}
Levy and Rousset gave the first complexity results for verifying rule
bases~\cite{LevyRousset1996,LevyRousset1998}: whether no legal input derives a bad output, for
relational inputs and Horn rules, NP-complete to undecidable depending on recursion, and
$\Pi_2^p$-complete with order predicates and negation. Their ``input completeness'' constraints
restrict which inputs are legal, and they leave output completeness open. Coverage of a
propositional table could be phrased in their framework, but its complexity is lower: coNP for
completeness, as shown here, with NP for non-emptiness and \#P for counting.

To our knowledge, none of these works classifies the completeness of decision tables with complete
tuples as rows over declared finite domains, or counts their missing combinations, and none of the
hardness proofs is machine-checked.

\subparagraph*{Mechanised computability and complexity.}
Cook--Levin has been mechanised in ACL2~\cite{GamboaCowles2004}, Coq~\cite{GaeherKunze2021} and
Isabelle~\cite{Cook_Levin-AFP}; reductions between Karp's problems are verified in
Isabelle~\cite{PolyReductionsIsabelle}; and Senellart and Gnatenko prove all 21 of Karp's
problems NP-complete in Lean~\cite{SenellartGnatenko2026}, with NP defined in
descriptive-complexity terms and first-order reductions, and also prove \#Monotone-2SAT one-call
\#P-complete in that setting. Turing machines and computability have been mechanised in
Isabelle/HOL~\cite{XuZhangUrban2013} and Coq~\cite{ForsterKunzeWuttke2020}, computability through
partial recursive functions in Lean~\cite{Carneiro2019}, and time complexity through the weak
call-by-value $\lambda$-calculus in Coq~\cite{ForsterKSW2021}. Our setting needs explicit encodings and running
times, but its classes are the textbook ones on Turing machines, so polynomial time, binary
encodings of numbers and the size of the input, the quantities the tractability results of
\Cref{sec:boundary} are about, are first-class.

%-----------------------------------------------------------------------------------------------
\section{The formalisation}
\label{sec:formal}

\subparagraph*{Trust.}
The axiom audit prints the axioms of 318 of the development's 335 theorem and lemma declarations,
among them every theorem this paper cites (the other 17 are helper lemmas used by audited
theorems); each uses only \lstinline{propext}, \lstinline{Classical.choice} and \lstinline{Quot.sound}, or
fewer, and none uses \lstinline{sorryAx}. This shows that no proof is left unfinished and no
axiom is added. It does not show that the definitions are the intended ones. The trusted base is
Lean's kernel; Complexitylib's definitions of machines, running time, encodings and the classes
P, NP, coNP and \#P; our encodings of instances; \Cref{def:onecall}; and the problem definitions
themselves: \textit{MissingNonempty}, \textit{missingCount}, the monotone shape tests and
\textit{SharpMonotone2SAT}. The adequacy theorems shrink the list: the general problems are
proved equal to latc-check's, so for them a reader need trust only latc-check's definition of the
missing set and our map from an instance to the table it denotes, each a few lines
(\Cref{sec:prelim,sec:encoding}). For the monotone counts, adequacy holds only on instances that
pass the shape test, so the shape test is trusted; an earlier version of it accepted conditions on
nonexistent attributes, which the present one rejects.

\subparagraph*{Engineering.}
The development has 4\,418 lines of Lean, besides the audit files. Mathlib's
\lstinline{fun_prop} tactic, with Complexitylib's polynomial-time lemmas, discharged most composition steps of polynomial-time proofs. Proofs by hand
were needed where the size of intermediate values matters: the binary product of the domain
sizes, the enumeration of candidate states (by induction on~$k$), the explicit bounds of
\Cref{sec:boundary}, and the arithmetic of Valiant's post-processing. For the graph of
\Cref{sec:monotone}, \lstinline{fun_prop} timed out until each vertex-numbering and edge helper
had its own composable polynomial-time lemma. Reading a list at a binary index needed a library
lemma, \lstinline{listGetD_binary}, rather than a new proof. The development pins one commit of
Complexitylib, a library that changes daily on a release-candidate toolchain; a version of the
development taken before the last two monotone changes also built unchanged against a later
commit on the next toolchain (\Cref{app:artifact}).

\subparagraph*{Worst case and practice.}
The classification rules out a polynomial algorithm for all tables with clause constraints,
unless P~=~NP; it does not say that tables met in practice are slow to check. For tables whose
rows are complete tuples, few input columns or the absence of constraints give the polynomial
regimes of \Cref{sec:boundary}, and the difficulty comes from constraints across many
attributes, not from the size of the domains or of the table. For compact rows, as in DMN, the
rows themselves can carry the hardness (\Cref{sec:related}), and our tractability results do not
apply.

%-----------------------------------------------------------------------------------------------
\section{Conclusion}
\label{sec:conclusion}

For decision tables whose rows are complete tuples, completeness under clause constraints is
coNP-complete, finding a missing combination NP-complete, and counting them \#P-complete, with
machine-checked proofs on Turing machines that apply to latc-check's set-based definitions. All
three are polynomial without constraints; they are also polynomial with a fixed number of
attributes, with a degree linear in that number. For monotone two-literal constraints, counting
the missing combinations stays \#P-complete under one-call reductions, on tables with arbitrary
rows and already with none, by a one-call version of Valiant's theorem proved in the same model. What remains
open in the formalisation is the per-$k$ machine with constants independent of~$k$. The one-call
reductions and the Valiant reduction use only Complexitylib's notions and could be offered
upstream.

%-----------------------------------------------------------------------------------------------
\bibliography{../docs/paper-refs,extra}

\appendix

%-----------------------------------------------------------------------------------------------
\section{Main Lean statements}
\label{app:statements}

The statements below are copied from the source up to \lstinline{:=}. Unless noted, they are in
the namespace \lstinline{Latc.PaperD.General}; \lstinline{Inst} is the type of instances of
\Cref{sec:encoding}.

\subparagraph*{Bridge (namespace \lstinline{Latc.PaperD}).}
\begin{lstlisting}
theorem holdsAll_eq_eval (S : Prop5.ConSet ℕ) (x : ℕ → Bool) :
    Prop5.holdsAll S x = (toCNF S).eval x
theorem satisfiable_decode_iff (S : Prop5.ConSet ℕ) :
    Prop5.Satisfiable (Prop5.decode S) ↔ (toCNF S).Satisfiable
\end{lstlisting}

\subparagraph*{Instances.}
\begin{lstlisting}
abbrev Cond := ℕ × ℕ
abbrev Inst := List ℕ × (List (List Cond) × List (List ℕ))
def Valid (S : List (List Cond)) (s : List ℕ) : Prop := ∀ c ∈ S, ∃ a ∈ c, CondHolds s a
def MissingNonempty : Set Inst :=
  {I | ∃ s, InDomains I.1 s ∧ Valid I.2.1 s ∧ s ∉ I.2.2}
def Complete : Set Inst := MissingNonemptyᶜ
\end{lstlisting}

\subparagraph*{The Boolean model (namespace \lstinline{Latc.PaperD}).}
\begin{lstlisting}
theorem MissingNonempty_npComplete : NPComplete encInst MissingNonempty
theorem Complete_coNPComplete : coNPComplete encInst Complete
theorem countMissing_sharpPComplete : SharpPComplete encInst countMissing
\end{lstlisting}
Here \lstinline{Inst}, \lstinline{encInst} and \lstinline{MissingNonempty} are those of the
Boolean model, a different namespace from the general one below.

\subparagraph*{Decision problems and counting.}
\begin{lstlisting}
theorem MissingNonempty_npComplete_bin : NPComplete encInstBin MissingNonempty
theorem Complete_coNPComplete_bin : coNPComplete encInstBin Complete
theorem missingNonempty_finset_npComplete_bin :
    NPComplete encInstBin {I : Inst | (missing I.con I.table).Nonempty}
theorem complete_finset_coNPComplete_bin :
    coNPComplete encInstBin {I : Inst | missing I.con I.table = ∅}
theorem missingCount_eq_acceptCount (I : Inst) :
    missingCount I = acceptCount cverifier I (certLen I)
theorem sharpSAT_eq_missingCount (φ : CNF ℕ) : SharpSAT φ = missingCount (fromSAT φ)
theorem missingCount_sharpPComplete_bin : SharpPComplete encInstBin missingCount
\end{lstlisting}

\subparagraph*{One-call reductions and Valiant.}
\lstinline{OneCallReducible} is in \lstinline{Latc.PaperD.General}; the four theorems after it are in
\lstinline{Latc.PaperD.Valiant}.
\begin{lstlisting}
def OneCallReducible (encA : α ↪ List Bool) (encB : β ↪ List Bool) (f : α → ℕ) (g : β → ℕ) :
    Prop :=
  ∃ (r : α → β) (post : α → ℕ → ℕ), ComputableInPolyTime r encA encB ∧
    ComputableInPolyTime (fun p : α × ℕ => post p.1 p.2) (encPair encA encNat) encNat ∧
      ∀ a, f a = post a (g (r a))
theorem sharpMonotone2SAT_oneCallSharpPHard :
    OneCallSharpPHard encCNF SharpMonotone2SAT
theorem countValidMonotone2_sharpPComplete_oneCall :
    OneCallSharpPComplete encInstBin countValidMonotone2
theorem countMissingMonotone2_eq_card_missing (I : General.Inst) (h : pos2AnyB I = true) :
    countMissingMonotone2 I = (Latc.missing I.con I.table).card
theorem countMissingMonotone2_sharpPComplete_oneCall :
    OneCallSharpPComplete encInstBin countMissingMonotone2
\end{lstlisting}

\subparagraph*{Tractability.}
\begin{lstlisting}
theorem missingCount_eq_of_noConstraints (I : Inst) (h : I.2.1 = []) :
    missingCount I = domProd I - distinctRows I
theorem missingCountFixedAttrs_polyTime (k : ℕ) :
    ComputableInPolyTime (fun I : Inst => if I.n ≤ k then missingCount I else 0)
      encInstBin encNat
theorem missingNonempty_explicitDegree :
    ∃ c d, ComputableTidilyInTimeAndSpace
      (fun I : Inst => decide (missing I.con I.table).Nonempty) encInstBin encBool
      (fun I => c * ((encInstBin I).length + 1) ^ (d * I.n + d))
      (fun I => c * ((encInstBin I).length + 1) ^ (d * I.n + d))
theorem missingNonemptyFixedAttrs_uniform :
    ∃ c d₀ d₁, ∀ k, ComputableTidilyInTimeAndSpace
      (fun I : {I : Inst // I.n ≤ k} => decide (missing I.1.con I.1.table).Nonempty)
      ((Function.Embedding.subtype _).trans encInstBin) encBool
      (fun I => c * ((encInstBin I.1).length + 1) ^ (d₀ * k + d₁))
      (fun I => c * ((encInstBin I.1).length + 1) ^ (d₀ * k + d₁))
\end{lstlisting}

%-----------------------------------------------------------------------------------------------
\section{Artifact and reproducibility}
\label{app:artifact}

\subparagraph*{Contents.}
The development is the directory \texttt{LatcComplexity/} of the repository latc-complexity:
4\,788 lines of Lean in 16 files, 4\,418 without the two axiom-audit files, on top of
latc-check's model, copied unchanged into \texttt{LatcFormal/}. The README lists every theorem
with its statement, and \texttt{docs/paper-facts.md} records the file and line of every theorem
this paper cites.

\subparagraph*{Versions.}
Lean \texttt{v4.36.0-rc1}; Complexitylib at commit \texttt{b7acfc2} (28 October 2026), with the
CSLib and Mathlib it pins. The pin was kept throughout the development. Complexitylib had just
been rewritten on CSLib; the earlier version, which Russell~\cite{PaperA} cites, is a separate
repository with a different API.

\subparagraph*{Build and audit.}
From a fresh clone, downloading precompiled Mathlib, \texttt{lake build} completed 18\,645 jobs
with no warnings or errors in about nine minutes on four cores, using 12\,GB of disk, at the
commit before the last theorem file was added; with that file, and with the monotone shape tests
tightened, the build has 18\,647 jobs, again clean. The axiom audit (the two files \texttt{Axioms.lean} of \texttt{LatcComplexity/} and of
its subdirectory \texttt{Valiant/}) prints 318 theorems: 268 use all three standard axioms, 50
fewer, none \lstinline{sorryAx}. The development as of commit \texttt{6603497}, before the
monotone result for arbitrary rows and the tightened shape tests, also built without any change
against Complexitylib commit \texttt{b9e6241} on Lean \texttt{v4.36.0-rc2} (18\,717 jobs); the
final version was not re-tested against it. The artifact is published with the pin above,
against which the results were proved.

\subparagraph*{History.}
The core results were committed in seven commits between 28 October 2026, 14:16, and
29 October, 00:18. On 31 October, counting with a fixed number of attributes, the Valiant
reduction (1\,140 lines when merged, besides its audit file) and
the explicit degree (692 lines) were added, the last two in parallel and merged without
conflicts; then the monotone result for tables with any rows, and a check in the monotone shape
tests that every condition names an existing attribute.
% Upstream figures from cslib-community/complexitylib, main at ff6ac8e on 2026-10-31; PRs #1653
% (Lean rc2), #1694 (CSLib and Mathlib on rc2), #1657 (fun_prop clients on rc2). Upstream test:
% branch claude/upstream-complexitylib-test, Complexitylib b9e6241, 18717 jobs, clean.

\end{document}
